%% BioMed_Central_Tex_Template_v1.06
%%                                      %
%  bmc_article.tex            ver: 1.06 %
%                                       %

%%IMPORTANT: do not delete the first line of this template
%%It must be present to enable the BMC Submission system to
%%recognise this template!!

%%%%%%%%%%%%%%%%%%%%%%%%%%%%%%%%%%%%%%%%%
%%                                     %%
%%  LaTeX template for BioMed Central  %%
%%     journal article submissions     %%
%%                                     %%
%%          <8 June 2012>              %%
%%                                     %%
%%                                     %%
%%%%%%%%%%%%%%%%%%%%%%%%%%%%%%%%%%%%%%%%%


%%%%%%%%%%%%%%%%%%%%%%%%%%%%%%%%%%%%%%%%%%%%%%%%%%%%%%%%%%%%%%%%%%%%%
%%                                                                 %%
%% For instructions on how to fill out this Tex template           %%
%% document please refer to Readme.html and the instructions for   %%
%% authors page on the biomed central website                      %%
%% http://www.biomedcentral.com/info/authors/                      %%
%%                                                                 %%
%% Please do not use \input{...} to include other tex files.       %%
%% Submit your LaTeX manuscript as one .tex document.              %%
%%                                                                 %%
%% All additional figures and files should be attached             %%
%% separately and not embedded in the \TeX\ document itself.       %%
%%                                                                 %%
%% BioMed Central currently use the MikTex distribution of         %%
%% TeX for Windows) of TeX and LaTeX.  This is available from      %%
%% http://www.miktex.org                                           %%
%%                                                                 %%
%%%%%%%%%%%%%%%%%%%%%%%%%%%%%%%%%%%%%%%%%%%%%%%%%%%%%%%%%%%%%%%%%%%%%

%%% additional documentclass options:
%  [doublespacing]
%  [linenumbers]   - put the line numbers on margins

%%% loading packages, author definitions

%\documentclass[twocolumn]{bmcart}% uncomment this for twocolumn layout and comment line below
\documentclass{bmcart}

%%% Load packages
%\usepackage{amsthm,amsmath}
%\RequirePackage{natbib}
%\RequirePackage[authoryear]{natbib}% uncomment this for author-year bibliography
%\RequirePackage{hyperref}
\usepackage[utf8]{inputenc} %unicode support
%\usepackage[applemac]{inputenc} %applemac support if unicode package fails
%\usepackage[latin1]{inputenc} %UNIX support if unicode package fails


%%%%%%%%%%%%%%%%%%%%%%%%%%%%%%%%%%%%%%%%%%%%%%%%%
%%                                             %%
%%  If you wish to display your graphics for   %%
%%  your own use using includegraphic or       %%
%%  includegraphics, then comment out the      %%
%%  following two lines of code.               %%
%%  NB: These line *must* be included when     %%
%%  submitting to BMC.                         %%
%%  All figure files must be submitted as      %%
%%  separate graphics through the BMC          %%
%%  submission process, not included in the    %%
%%  submitted article.                         %%
%%                                             %%
%%%%%%%%%%%%%%%%%%%%%%%%%%%%%%%%%%%%%%%%%%%%%%%%%


%\def\includegraphic{}
%\def\includegraphics{}


%%TOMMI
%ricordati di scommentare le righe sopra!
\usepackage{graphicx, subfigure}

\newcommand{\ie}{\textit{i.e.}~}
\newcommand{\iec}{\textit{i.e.},~}
\newcommand{\eg}{\textit{e.g.}~}
\newcommand{\egc}{\textit{e.g.},~}

\newcommand{\eal}{\textit{et al.}~}
\newcommand{\cfr}{cfr.~}

\newcommand{\fig}{Fig.~}
\newcommand{\figs}{Figs.~}




%%% Put your definitions there:
\startlocaldefs
\endlocaldefs


%%% Begin ...
\begin{document}

%%% Start of article front matter
\begin{frontmatter}

\begin{fmbox}
\dochead{Research}

%%%%%%%%%%%%%%%%%%%%%%%%%%%%%%%%%%%%%%%%%%%%%%
%%                                          %%
%% Enter the title of your article here     %%
%%                                          %%
%%%%%%%%%%%%%%%%%%%%%%%%%%%%%%%%%%%%%%%%%%%%%%

\title{Design and evaluation of a low-cost instrumented glove for thumb rehabilitation and functionality assessment}

%%%%%%%%%%%%%%%%%%%%%%%%%%%%%%%%%%%%%%%%%%%%%%
%%                                          %%
%% Enter the authors here                   %%
%%                                          %%
%% Specify information, if available,       %%
%% in the form:                             %%
%%   <key>={<id1>,<id2>}                    %%
%%   <key>=                                 %%
%% Comment or delete the keys which are     %%
%% not used. Repeat \author command as much %%
%% as required.                             %%
%%                                          %%
%%%%%%%%%%%%%%%%%%%%%%%%%%%%%%%%%%%%%%%%%%%%%%

\author[
   addressref={unisi},                   % id's of addresses, e.g. {aff1,aff2}
   corref={unisi},                       % id of corresponding address, if any
   %noteref={n1},                        % id's of article notes, if any
   email={pompili@diism.unisi.it}   % email address
]{\inits{GP}\fnm{Giulia} \snm{Pompili}}
%
\author[
   addressref={unisi,iit},%
   email={lisini@diism.unisi.it}%
]{\inits{TLB}\fnm{Tommaso} \snm{Lisini Baldi}}
%
\author[
   addressref={unisi,iit},
   email={scal@diism.unisi.it}
]{\inits{SC}\fnm{Samanta} \snm{Scala}}
%
\author[
   addressref={unisi},
   email={meli@diism.unisi.it}
]{\inits{LM}\fnm{Leonardo} \snm{Meli}}
%
\author[
   addressref={unisi,iit},
   email={prattichizzo@diism.unisi.it}
]{\inits{DP}\fnm{Domenico} \snm{Prattichizzo}}
%%%%%%%%%%%%%%%%%%%%%%%%%%%%%%%%%%%%%%%%%%%%%%
%%                                          %%
%% Enter the authors' addresses here        %%
%%                                          %%
%% Repeat \address commands as much as      %%
%% required.                                %%
%%                                          %%
%%%%%%%%%%%%%%%%%%%%%%%%%%%%%%%%%%%%%%%%%%%%%%

\address[id=unisi]{%                           % unique id
  \orgname{Department of Information Engineering and Mathematics, University of Siena}, % university, etc
 \street{via Roma, 56},                     %
 \postcode{53100}                                % post or zip code
 \city{Siena},                              % city
 \cny{Italy}                                 % country
}
  
\address[id=iit]{%
\orgname{Department of Advanced Robotics, Istituto Italiano di Tecnologia},
 \street{Via Morego, 30},
  \postcode{16163}
 \city{Genova},
 \cny{Italy}
}

%%%%%%%%%%%%%%%%%%%%%%%%%%%%%%%%%%%%%%%%%%%%%%
%%                                          %%
%% Enter short notes here                   %%
%%                                          %%
%% Short notes will be after addresses      %%
%% on first page.                           %%
%%                                          %%
%%%%%%%%%%%%%%%%%%%%%%%%%%%%%%%%%%%%%%%%%%%%%%

\begin{artnotes}
%\note{Sample of title note}     % note to the article
%\note[id=aff1]{Equal contributor} % note, connected to author
\end{artnotes}

\end{fmbox}% comment this for two column layout

%%%%%%%%%%%%%%%%%%%%%%%%%%%%%%%%%%%%%%%%%%%%%%
%%                                          %%
%% The Abstract begins here                 %%
%%                                          %%
%% Please refer to the Instructions for     %%
%% authors on http://www.biomedcentral.com  %%
%% and include the section headings         %%
%% accordingly for your article type.       %%
%%                                          %%
%%%%%%%%%%%%%%%%%%%%%%%%%%%%%%%%%%%%%%%%%%%%%%

\begin{abstractbox}
%
\begin{abstract} % abstract
\parttitle{First part title} %if any
Text for this section.
%
\parttitle{Second part title} %if any
Text for this section.
\end{abstract}
%
%%%%%%%%%%%%%%%%%%%%%%%%%%%%%%%%%%%%%%%%%%%%%%
%%                                          %%
%% The keywords begin here                  %%
%%                                          %%
%% Put each keyword in separate \kwd{}.     %%
%%                                          %%
%%%%%%%%%%%%%%%%%%%%%%%%%%%%%%%%%%%%%%%%%%%%%%
%
\begin{keyword}
\kwd{sample}
\kwd{article}
\kwd{author}
\end{keyword}
%
% MSC classifications codes, if any
%\begin{keyword}[class=AMS]
%\kwd[Primary ]{}
%\kwd{}
%\kwd[; secondary ]{}
%\end{keyword}
%
\end{abstractbox}
%
%\end{fmbox}% uncomment this for twcolumn layout
%
\end{frontmatter}

%%%%%%%%%%%%%%%%%%%%%%%%%%%%%%%%%%%%%%%%%%%%%%
%%                                          %%
%% The Main Body begins here                %%
%%                                          %%
%% Please refer to the instructions for     %%
%% authors on:                              %%
%% http://www.biomedcentral.com/info/authors%%
%% and include the section headings         %%
%% accordingly for your article type.       %%
%%                                          %%
%% See the Results and Discussion section   %%
%% for details on how to create sub-sections%%
%%                                          %%
%% use \cite{...} to cite references        %%
%%  \cite{koon} and                         %%
%%  \cite{oreg,khar,zvai,xjon,schn,pond}    %%
%%  \nocite{smith,marg,hunn,advi,koha,mouse}%%
%%                                          %%
%%%%%%%%%%%%%%%%%%%%%%%%%%%%%%%%%%%%%%%%%%%%%%

%%%%%%%%%%%%%%%%%%%%%%%%% start of article main body
% <put your article body there>

%%%%%%%%%%%%%%%%
%% Background %%
%%%%%%%%%%%%%%%%
\section*{Introduction}

The opposite thumb allows humans to distinguish themselves from the others animal species. The morphological hand configuration and the kinematic coordination between all fingers have a crucial role in the manipulative abilities, precision-pinch grips as well as power “squeeze” grips \cite{napier1956prehensile} \cite{marzke1997precision} \cite{marzke1992evolution}. Thumb, in particular, is the finger who has the larger representative area in the brain. It covers about 50\% of the interior cortical area related to the hand and its loss or absence causes the loss of 40\% of hand functionality \cite{flatt2002our}. 

Surgery is usually recommended for many injuries or degenerative diseases of this digit. Hand rehabilitation has to be focus to avoid or prevent stiffness caused by a prolonged immobilization after surgery.
The term \textit stiff is commonly used to describe an increasing rigidity but stiffness of the hand is more properly a loss of the elastic configuration of the collagen fibers \cite{frank1984physiology}. Active motion can provide several benefits for the strength of healing tissues and they are well described in literature \cite{akeson1987effects} \cite{cyr1998controlled} \cite{gelberman1981influence}, including the realignment and orientation of collagen fibers. Moreover, muscle contraction resulting from active motion can stimulate the motor cortex and the restoration (relearning) of normal movement patterns


\section*{Methods}
Text for this section \ldots
\subsection*{motivazioni rehab}
Kapandji test 
\subsection*{The glove}
In this section we present the proposed aid for rehabilitation, describing both the hardware component and the ad-hoc app developed for training and evaluating the hand mobility.
\subsubsection*{parte hardware}
The hardware part consist of a lycra glove, equipped with conductive bands to detect the thumbs contacts, and a tiny microcontroller in charge of indentifying the correspondent phalanx and send this information to the paired smartphone.
As depicted in \fig\ref{FIG:glove} we sewed 14 conductive fabric patches into the glove. 	

\begin{figure}
\label{FIG:glove}
\includegraphics[width= 3 cm]{./fig/blank.jpg}
\caption{Foto del guanto, magari subfig davanti e dietro}
\end{figure}



\subsubsection*{parte software}

\subsection*{Results}
Text for this sub-heading \ldots
\subsubsection*{Sub-sub heading for section}
Text for this sub-sub-heading \ldots

\bibliographystyle{bmc-mathphys} % Style BST file (bmc-mathphys, vancouver, spbasic).
\bibliography{biblio}      % Bibliography file (usually '*.bib' )
\end{document}
